%% kapa: exam=2 round=28 year=2017 period=3 subject=compensation_law subject_ko=감정평가 및 보상법규 seq=1 number=1 points=40
%% source: Q-Net 감정평가사 기출문제 (cst00309) artlSeq=5201712 file=raw/1008/5201712/2017년 감정평가사 2차시험 기출문제(감정평가 및 보상법규).pdf
%% conversion: kordoc 4.24.51 -> tools/gen2.py (auto); review=proofread-sonnet; flags=none
\begin{longquestion}{【문제 1】}
甲은 A시의 관할구역 내 X토지를 소유하고 있다. A시는 그동안 조선업의 지속적인 발전으로 다수의 인구가 거주하였으나 최근 세계적인 불황으로 인구가 급격하게 감소하고 있다. 국토교통부장관은 A시를 국제관광특구로 발전시킬 목적으로 「기업도시개발 특별법」이 정하는 바에 따라 X토지가 포함된 일단의 토지를 기업도시개발구역으로 지정하고, 개발사업시행자인 乙이 작성한 기업도시개발계획(동법 제14조 제2항에 따른 X토지 그 밖의 수용 대상이 되는 토지의 세부목록 포함. 이하 같다)을 승인·고시하였다. 乙은 협의취득에 관한 제반 절차를 준수하여 X토지에 대한 협의취득을 시도하였으나 甲이 응하지 않았다. 이에 乙은 X토지에 대한 수용재결을 신청하였고 중앙토지수용위원회는 그 신청에 따른 수용재결을 하였다. 다음 물음에 답하시오. (40점) \par
\medskip\noindent 1) 甲은 기업도시개발계획승인에 대한 취소소송의 제소기간이 도과한 상태에서, 「공익사업을 위한 토지 등의 취득 및 보상에 관한 법률」 제21조 제2항에 따른 중앙토지수용위원회 및 이해관계자의 의견청취절차를 전혀 시행하지 않은 채 기업도시개발계획승인이 발급된 것이 위법함을 이유로 수용재결 취소소송을 제기하려고 한다. 甲의 소송상 청구가 인용될 수 있는 가능성에 관하여 설명하시오. (단, 소송요건은 충족된 것으로 본다.) (20점) \par
\medskip\noindent 2) 甲은 수용재결 취소소송을 제기하면서, 乙이 기업도시개발계획승인 이후에 재정상황이 악화되어 수용재결 당시에 이르러 기업도시개발사업을 수행할 능력을 상실한 상태가 되었음에도 불구하고 수용재결을 한 위법이 있다고 주장한다. 甲의 소송상 청구가 인용될 수 있는 가능성에 관하여 설명하시오. (단, 소송요건은 충족된 것으로 본다.) (10점) \par
\medskip\noindent 3) 중앙토지수용위원회는 보상금을 산정하면서, X토지는 그 용도지역이 제1종 일반주거지역이기는 하지만 기업도시개발사업의 시행을 위해서 제3종 일반주거지역으로 변경되지 않은 사정이 인정되므로 제3종 일반주거지역으로 변경이 이루어진 상태를 상정하여 토지가격을 평가한다고 설시하였다. 이에 대해 乙은 X토지를 제1종 일반주거지역이 아닌 제3종 일반주거지역으로 평가한 것은 공법상 제한을 받는 토지에 대한 보상금 산정에 위법이 있다고 주장하면서 보상금감액청구소송을 제기하고자 한다. 乙의 소송상 청구가 인용될 수 있는 가능성에 관하여 설명하시오. (단, 소송요건은 충족된 것으로 본다.) (10점) \par
\medskip\noindent
\begin{framedtext}
\textbf{참조 조문} \par
\textbf{<공익사업을 위한 토지 등의 취득 및 보상에 관한 법률>} \par
\textbf{제21조(의견청취 등)} ① 생략 \par
② 별표에 규정된 법률에 따라 사업인정이 있는 것으로 의제되는 공익사업의 허가·인가·승인권자 등은 사업인정이 의제되는 지구지정·사업계획승인 등을 하려는 경우 제1항에 따라 제49조에 따른 중앙토지수용위원회 및 사업인정에 이해관계가 있는 자의 의견을 들어야 한다. \par
③ 생략 \par
\textbf{<공익사업을 위한 토지 등의 취득 및 보상에 관한 법률 시행규칙>} \par
\textbf{제23조(공법상 제한을 받는 토지의 평가)} ① 공법상 제한을 받는 토지에 대하여는 제한 받는 상태대로 평가한다. 다만, 그 공법상 제한이 당해 공익사업의 시행을 직접 목적으로 하여 가하여진 경우에는 제한이 없는 상태를 상정하여 평가한다. \par
② 당해 공익사업의 시행을 직접 목적으로 하여 용도지역 또는 용도지구 등이 변경된 토지에 대하여는 변경되기 전의 용도지역 또는 용도지구 등을 기준으로 평가한다. \par
\textbf{<기업도시개발 특별법>} \par
\textbf{제11조(개발계획의 승인 등)} ① 제4조에 따라 개발구역의 지정을 제안하는 자는 지정 제안 시 기업도시개발계획(이하 “개발계획”이라 한다)을 작성하여 국토교통부장관의 승인을 받아야 한다. <이하 생략> \par
②\textasciitilde{}③ 생략 \par
④ 국토교통부장관은 제1항에 따라 개발계획을 승인하였을 때에는 대통령령으로 정하는 바에 따라 그 사실을 관보에 게시하여야 한다. <이하 생략> \par
\textbf{제14조(토지등의 수용·사용)} ① 시행자는 개발구역에서 개발사업을 시행하기 위하여 필요할 때에는 「공익사업을 위한 토지 등의 취득 및 보상에 관한 법률」 제3조에 따른 토지·물건 또는 권리(이하 “토지등”이라 한다)를 수용 또는 사용(이하 “수용등”이라 한다)할 수 있다. \par
② 제1항을 적용하는 경우에 수용등의 대상이 되는 토지등의 세부목록을 제11조 제4항에 따라 고시한 때에는 「공익사업을 위한 토지 등의 취득 및 보상에 관한 법률」 제20조 제1항 및 제22조에 따른 사업인정 및 사업인정의 고시가 있은 것으로 본다. \par
③\textasciitilde{}⑨ 생략 \par
⑩ 제1항에 따른 토지등의 수용등에 관하여 이 법에 특별한 규정이 있는 경우를 제외하고는 「공익사업을 위한 토지 등의 취득 및 보상에 관한 법률」을 준용한다.
\end{framedtext}
\end{longquestion}
